%% ocean_fcl_routing.tex
%% Ocean FCL Multi-Commodity Routing — MVP Formal Model
%% Phase 1 — approved before any code starts
%% Status: Draft 2026-05-10

\documentclass[11pt,a4paper]{article}

\usepackage{amsmath,amssymb,amsthm}
\usepackage{booktabs}
\usepackage{geometry}
\usepackage{hyperref}
\usepackage{array}
\usepackage{xcolor}
\usepackage{enumitem}
\usepackage{longtable}
\usepackage{caption}

\geometry{margin=2.5cm}
\setlength{\parskip}{0.4em}

\hypersetup{colorlinks=true, linkcolor=blue, urlcolor=blue}

\newcommand{\ceil}[1]{\left\lceil #1 \right\rceil}
\newcommand{\defeq}{\triangleq}
\newcommand{\todo}[1]{\textcolor{red}{\textbf{[TODO: #1]}}}
\newcommand{\deferred}[1]{\textcolor{gray}{\textit{(Deferred P1: #1)}}}

\theoremstyle{definition}
\newtheorem{defn}{Definition}
\newtheorem{remark}{Remark}

\title{\textbf{Ocean FCL Multi-Commodity Routing}\\[0.5em]
\large MVP Formal Model --- Phase 1}
\author{AI Freight Agent}
\date{Draft --- 2026-05-10}

\begin{document}

\maketitle

\begin{abstract}
This document formally specifies the optimization model for the ocean FCL multi-commodity
routing problem, constituting the first component in the AI Freight Agent build sequence.
The model is a multi-commodity minimum-cost flow problem on a directed graph whose
structure is defined by the physical freight network (see \texttt{docs/ocean\_fcl\_multi\_shipment\_graph.drawio}).
Scope is limited to FCL (Full Container Load) shipments in 40$'$ containers (FEU).
LCL consolidation, probabilistic objectives, and commodity-specific customs inspection
probability are explicitly deferred to P1.
\end{abstract}

\tableofcontents
\newpage

%%--------------------------------------------------------------------
\section{Problem Statement}
%%--------------------------------------------------------------------

Given a set of shipment requests (commodities), each defined by an origin location,
destination location, cargo volume, cargo weight, time constraints, and service
attributes, find a minimum-cost assignment of commodities to paths through the
multimodal freight network such that:
\begin{itemize}[noitemsep]
  \item every commodity is delivered on or before its deadline,
  \item no ocean sailing exceeds its allocated FEU slot capacity, and
  \item optional per-commodity budget caps are respected.
\end{itemize}

The routing network consists of four arc types: pre-carriage truck (origin door to
port of loading), ocean sailing (port of loading to port of discharge arrival), port
dwell (discharge arrival to port exit, covering unloading and customs clearance), and
inland truck or intermodal rail (port exit to destination door).

\paragraph{MVP scope.}
\begin{itemize}[noitemsep]
  \item FCL only. LCL deferred (PRD Section 11, Open Question 8).
  \item Container type: 40$'$ standard (FEU, $\approx 67$\,CBM usable). TEU deferred to P1.
  \item Port dwell: fixed port-specific mean. Commodity-specific inspection model deferred to P1.
  \item Objective: minimize total freight cost subject to \emph{deterministic} time-window feasibility.
        Probabilistic objectives (chance constraints on P(on-time)) deferred to P1.
  \item No transshipment arcs in MVP. Single-leg ocean sailing only.
\end{itemize}

%%--------------------------------------------------------------------
\section{Graph Structure}
%%--------------------------------------------------------------------

The freight network is a directed graph $G = (N, A)$.

\subsection{Node Sets}

\[
  N = N_O \;\cup\; N_{\text{POL}} \;\cup\; N_{\text{POD},A} \;\cup\; N_{\text{POD},E} \;\cup\; N_D
\]

\begin{description}[style=nextline, leftmargin=3cm, labelwidth=2.8cm]
  \item[$N_O$] \textbf{Origin nodes.} One per shipper origin location (factory or warehouse
        door). Identified by geographic coordinates (lat, lon). Commodity source nodes.
        Example: SUZ (Suzhou factory), GZH (Guangzhou warehouse), SZN (Shenzhen facility).

  \item[$N_{\text{POL}}$] \textbf{Port of Loading (POL) nodes.} Container terminals at the
        origin country where shipments are gated in and loaded onto vessels. Example: SHA
        (Shanghai), NGB (Ningbo), SZX (Shenzhen/Yantian).

  \item[$N_{\text{POD},A}$] \textbf{Port of Discharge Arrival nodes.} One arrival node per
        destination port (e.g., $\text{LAX}_A$, $\text{LGB}_A$, $\text{SEA}_A$). The
        commodity arrives from the ocean leg. It cannot depart until the dwell arc is traversed.

  \item[$N_{\text{POD},E}$] \textbf{Port of Discharge Exit nodes.} One exit node per
        destination port (e.g., $\text{LAX}_E$, $\text{LGB}_E$, $\text{SEA}_E$). The
        commodity has cleared customs and is available for inland pickup.

  \item[$N_D$] \textbf{Destination nodes.} One per consignee delivery location (warehouse
        or DC door). Commodity sink nodes.
        Example: NYC (New York DC), CHI (Chicago warehouse), DAL (Dallas DC).
\end{description}

\begin{remark}[POD node expansion]
Each physical port of discharge is represented by \emph{two} nodes connected by a
dwell arc: an arrival node $p_A \in N_{\text{POD},A}$ and an exit node
$p_E \in N_{\text{POD},E}$. This separation makes port dwell an explicit model
element rather than an implicit constant baked into inland arc transit times.
It also provides the structural hook for the P1 commodity-specific customs
inspection model (Section~\ref{sec:deferred}), where dwell arc weights will
become functions of commodity and shipper attributes rather than port-level constants.
\end{remark}

\subsection{Arc Sets}

\[
  A = A_{\text{PC}} \;\cup\; A_{\text{OC}} \;\cup\; A_{\text{DW}} \;\cup\; A_{\text{IL}}
\]

\begin{description}[style=nextline, leftmargin=3cm, labelwidth=2.8cm]
  \item[$A_{\text{PC}}$] \textbf{Pre-carriage arcs.} $(i, j)$ with $i \in N_O$,
        $j \in N_{\text{POL}}$. Truck move from origin door to container terminal.
        One arc per feasible (origin, POL) pair.

  \item[$A_{\text{OC}}$] \textbf{Ocean sailing arcs.} $(j, p)$ with $j \in N_{\text{POL}}$,
        $p \in N_{\text{POD},A}$. One arc per scheduled carrier sailing (specific vessel,
        departure date, and trade lane). Multiple sailings may exist on the same lane.

  \item[$A_{\text{DW}}$] \textbf{Dwell arcs.} $(p_A, p_E)$ with $p_A \in N_{\text{POD},A}$,
        $p_E \in N_{\text{POD},E}$. One arc per port of discharge. Represents container
        discharge from vessel plus customs clearance processing.

  \item[$A_{\text{IL}}$] \textbf{Inland arcs.} $(p_E, h)$ with $p_E \in N_{\text{POD},E}$,
        $h \in N_D$. Truck (FTL/drayage) or intermodal rail move from port exit to
        destination door.
\end{description}

%%--------------------------------------------------------------------
\section{Sets and Indices}
%%--------------------------------------------------------------------

\begin{center}
\begin{tabular}{ll}
\toprule
Symbol & Description \\
\midrule
$K$ & Set of commodities (shipment requests), indexed by $k$ \\
$N$ & Set of all nodes \\
$A$ & Set of all directed arcs \\
$A_{\text{PC}},\, A_{\text{OC}},\, A_{\text{DW}},\, A_{\text{IL}}$ & Arc subsets by type \\
$A^k \subseteq A$ & Feasible arc set for commodity $k$ (Section~\ref{sec:prefilter}) \\
$A^k_S \defeq A^k \cap A_S$ & Feasible arcs of type $S$ for commodity $k$ \\
\bottomrule
\end{tabular}
\end{center}

%%--------------------------------------------------------------------
\section{Parameters}
%%--------------------------------------------------------------------

\subsection{Commodity Parameters}

For each commodity $k \in K$:

\begin{center}
\begin{tabular}{llp{6cm}l}
\toprule
Parameter & Symbol & Description & Unit \\
\midrule
Volume & $v_k$ & Cargo volume & CBM \\
Weight & $w_k$ & Gross weight & kg \\
FEU count & $n_k = \ceil{v_k / 67}$ & Containers required; pre-computed, not a decision variable & FEU \\
Cargo ready & $t_k^{\text{rdy}}$ & Earliest pickup datetime at origin & datetime \\
Deadline & $T_k^{\text{dead}}$ & Latest acceptable arrival at destination & datetime \\
Origin & $o(k) \in N_O$ & Commodity source node & --- \\
Destination & $d(k) \in N_D$ & Commodity sink node & --- \\
Cargo type & $\tau_k$ & General / hazmat class / temp-controlled / OOG & categorical \\
HS code & $\text{hs}_k$ & 6-digit Harmonized System commodity code & --- \\
Incoterm & $\text{inco}_k$ & EXW / FOB / CIF / DDP (determines arc cost responsibility) & --- \\
Service level & $\text{sl}_k$ & Economy / Standard / Express & categorical \\
Budget cap & $B_k$ & Hard cost ceiling; $B_k = \infty$ if not set & USD \\
\bottomrule
\end{tabular}
\end{center}

\begin{remark}[HS code use]
The HS code $\text{hs}_k$ is an input in MVP but is not used in the MVP objective or
constraints --- it is captured for forward compatibility with the P1 customs inspection
model. The HS chapter (first two digits) will map to an inspection-risk tier.
\end{remark}

\begin{remark}[Incoterm and arc cost responsibility]
Under EXW, the buyer (importer) pays all costs from the origin door.
Under DDP, the seller (exporter) pays all costs to the destination door.
In MVP, the objective minimizes total door-to-door cost regardless of Incoterm.
Incoterm is recorded as an attribute and will be used in Phase 2 to split cost
responsibility and generate correct carrier invoices.
\end{remark}

\subsection{Pre-Carriage Arc Parameters}

For each $(i, j) \in A_{\text{PC}}$:

\begin{center}
\begin{tabular}{llp{6.5cm}l}
\toprule
Parameter & Symbol & Description & Unit \\
\midrule
Cost & $c_{ij}^{\text{PC}}$ & Flat truck rate per move. FCL economics: one chartered truck regardless of fill. Not per-CBM. & USD \\
Mean transit & $\mu_{ij}^{\text{PC}}$ & Mean time from origin door to terminal gate & days \\
Std dev transit & $\sigma_{ij}^{\text{PC}}$ & Std dev of transit time & days \\
Road distance & $\text{dist}_{ij}^{\text{PC}}$ & Road distance; informational, used by instance generator & km \\
\bottomrule
\end{tabular}
\end{center}

\subsection{Ocean Sailing Arc Parameters}

For each $(j, p) \in A_{\text{OC}}$:

\begin{center}
\begin{tabular}{llp{6cm}l}
\toprule
Parameter & Symbol & Description & Unit \\
\midrule
ETD & $\text{ETD}_{jp}$ & Estimated time of departure from POL $j$ & datetime \\
CY cutoff & $\text{CYC}_{jp}$ & Latest datetime a container may arrive at the terminal for this sailing. $\text{CYC}_{jp} \leq \text{ETD}_{jp}$, typically $\text{ETD} - 4\,\text{days}$. & datetime \\
Base rate & $r_{jp}^{\text{FEU}}$ & Ocean freight rate per 40$'$ container & USD/FEU \\
BAF & $\text{BAF}_{jp}$ & Bunker adjustment factor per FEU & USD/FEU \\
THC (POL) & $\text{THC}_{jp}^{\text{POL}}$ & Terminal handling charge at origin port per FEU & USD/FEU \\
THC (POD) & $\text{THC}_{jp}^{\text{POD}}$ & Terminal handling charge at destination port per FEU & USD/FEU \\
Capacity & $\text{cap}_{jp}$ & Maximum FEU slots available on this sailing (carrier-allocated block, not full vessel) & FEU \\
Mean transit & $\mu_{jp}^{\text{OC}}$ & Mean ocean transit time, POL $j$ to POD $p$ & days \\
Std dev transit & $\sigma_{jp}^{\text{OC}}$ & Std dev of ocean transit time & days \\
Carrier & $\text{carrier}_{jp}$ & Carrier identifier (e.g., MSC, CMA CGM, COSCO) & --- \\
Vessel & $\text{vessel}_{jp}$ & Vessel name & --- \\
\bottomrule
\end{tabular}
\end{center}

\noindent Derived all-in ocean cost per FEU:
\[
  c_{jp}^{\text{OC}} \;\defeq\; r_{jp}^{\text{FEU}} + \text{BAF}_{jp} + \text{THC}_{jp}^{\text{POL}} + \text{THC}_{jp}^{\text{POD}}
\]

\begin{remark}[CY cutoff semantics]
The CY (Container Yard) cutoff $\text{CYC}_{jp}$ is the hard deadline by which the
loaded, sealed container must be physically present and registered at the terminal for
the sailing. Missing the CY cutoff means the container is rolled to the next sailing.
A separate VGM (Verified Gross Mass) cutoff typically precedes CYC by 24--48 hours;
doc cutoff (SI/BL instructions) may coincide with CYC or fall 1--2 days earlier.
In MVP, $\text{CYC}_{jp}$ is the binding time constraint on the pre-carriage leg.
\end{remark}

\begin{remark}[POD arrival window]
The expected arrival of a commodity at $p_A$ is $\text{ETD}_{jp} + \mu_{jp}^{\text{OC}}$
(mean) with uncertainty $\sigma_{jp}^{\text{OC}}$. In MVP, the deterministic expected
arrival $\text{ETD}_{jp} + \mu_{jp}^{\text{OC}} + \mu_p^{\text{DW}}$ is used in the
feasibility pre-filter. In P1, a chance constraint will replace this with
$\Pr(\text{arrival} \leq T_k^{\text{dead}}) \geq \alpha_k$.
\end{remark}

\subsection{Dwell Arc Parameters}

For each $(p_A, p_E) \in A_{\text{DW}}$, parameterized by port $p$:

\begin{center}
\begin{tabular}{llp{5.5cm}l}
\toprule
Parameter & Symbol & Description & MVP value \\
\midrule
Mean dwell & $\mu_p^{\text{DW}}$ & Expected total dwell: unloading from vessel + customs clearance & USLAX: 3.5d; USLGB: 3.5d; USSEA: 2.5d \\
Std dev dwell & $\sigma_p^{\text{DW}}$ & Std dev of total dwell time & USLAX: 1.5d; USLGB: 1.5d; USSEA: 1.0d \\
Free time & $f_p$ & Port-provided free days before container storage charges apply & 5 days (standard) \\
Storage rate & $s_p$ & Per-diem cost after free time & Deferred to P1 \\
\bottomrule
\end{tabular}
\end{center}

\noindent Dwell time components:
\begin{align*}
  \mu_p^{\text{DW}} &= \mu_p^{\text{unload}} + \mu_p^{\text{clear}} \\
  \mu_p^{\text{unload}} &\approx 0.5\text{--}1.0\,\text{d} \quad \text{(discharge to accessible in terminal)} \\
  \mu_p^{\text{clear}} &\approx 1.0\text{--}3.0\,\text{d} \quad \text{(CBP processing, no exam; port-level mean)}
\end{align*}

\noindent \textit{P1 extension:} In P1, $\mu_p^{\text{DW}}$ becomes a function of commodity
attributes --- HS code, importer C-TPAT status, country of origin, and consignee history.
See Section~\ref{sec:deferred}.

\subsection{Inland Arc Parameters}

For each $(p_E, h) \in A_{\text{IL}}$:

\begin{center}
\begin{tabular}{llp{5.5cm}l}
\toprule
Parameter & Symbol & Description & Unit \\
\midrule
Cost & $c_{p_E h}^{\text{IL}}$ & Flat rate per container move (truck or rail) & USD/FEU \\
Mean transit & $\mu_{p_E h}^{\text{IL}}$ & Mean transit time, POD exit to destination door & days \\
Std dev transit & $\sigma_{p_E h}^{\text{IL}}$ & Std dev of transit time & days \\
Mode & $\text{mode}_{p_E h}$ & FTL truck or intermodal rail & categorical \\
Road distance & $\text{dist}_{p_E h}^{\text{IL}}$ & Road/rail distance; informational & km \\
\bottomrule
\end{tabular}
\end{center}

%%--------------------------------------------------------------------
\section{Feasibility Pre-Filtering}
\label{sec:prefilter}
%%--------------------------------------------------------------------

Before constructing the MILP, we compute the feasible arc set $A^k$ for each
commodity $k$. An arc is feasible for commodity $k$ only if it satisfies all time
constraints. This reduces the variable count and eliminates infeasible paths from
the formulation entirely, keeping the MILP small.

\paragraph{Pre-carriage feasibility.}
Arc $(i, j) \in A_{\text{PC}}$ is in $A^k$ only if:
\begin{enumerate}[noitemsep]
  \item $i = o(k)$ \quad (arc originates at commodity's origin node), and
  \item $\exists\, (j, p) \in A_{\text{OC}}$ such that $t_k^{\text{rdy}} + \mu_{ij}^{\text{PC}} \leq \text{CYC}_{jp}$. \\
        (Cargo can be delivered to terminal before the CY cutoff of at least one sailing.)
\end{enumerate}

\paragraph{Ocean arc compatibility.}
Ocean arc $(j, p) \in A_{\text{OC}}$ is compatible with pre-carriage arc
$(i, j) \in A^k_{\text{PC}}$ for commodity $k$ if:
\[
  t_k^{\text{rdy}} + \mu_{ij}^{\text{PC}} \leq \text{CYC}_{jp}
\]
Only compatible (pre-carriage, ocean) pairs generate connected paths in the graph.

\paragraph{End-to-end feasibility.}
The complete path (pre-carriage $(i,j)$, ocean $(j,p)$, dwell, inland $(p_E, h)$)
is feasible for commodity $k$ if:
\begin{enumerate}[noitemsep]
  \item $h = d(k)$ \quad (inland arc terminates at commodity's destination node), and
  \item $\text{ETD}_{jp} + \mu_{jp}^{\text{OC}} + \mu_p^{\text{DW}} + \mu_{p_E h}^{\text{IL}} \leq T_k^{\text{dead}}$.
\end{enumerate}

Dwell arcs $A_{\text{DW}}$ have no feasibility condition beyond arc connectivity; they are
always in $A^k$ if the associated POD appears in a feasible path.

%%--------------------------------------------------------------------
\section{Decision Variables}
%%--------------------------------------------------------------------

\[
  x_{ij}^k \in \{0, 1\} \qquad \forall\, k \in K,\; (i,j) \in A^k
\]

$x_{ij}^k = 1$ if commodity $k$ uses arc $(i,j)$; zero otherwise.

%%--------------------------------------------------------------------
\section{Objective Function}
%%--------------------------------------------------------------------

Minimize total freight cost across all commodities:

\begin{equation}
  \min \;\; Z = \sum_{k \in K} \Bigg[
    \underbrace{\sum_{(i,j) \in A^k_{\text{PC}}} c_{ij}^{\text{PC}}\, x_{ij}^k}_{\text{pre-carriage}}
    \;+\;
    \underbrace{\sum_{(j,p) \in A^k_{\text{OC}}} n_k\, c_{jp}^{\text{OC}}\, x_{jp}^k}_{\text{ocean (per FEU)}}
    \;+\;
    \underbrace{\sum_{(p_E,h) \in A^k_{\text{IL}}} n_k\, c_{p_E h}^{\text{IL}}\, x_{p_E h}^k}_{\text{inland (per FEU)}}
  \Bigg]
  \label{eq:obj}
\end{equation}

\paragraph{Cost structure notes.}
\begin{itemize}[noitemsep]
  \item \textit{Pre-carriage:} $c_{ij}^{\text{PC}}$ is a flat truck move rate. Under FCL
        economics, the shipper charters the truck; cost does not scale with CBM.
  \item \textit{Ocean:} Cost scales by $n_k$ (number of FEU) since each container is a
        separately quoted unit. $c_{jp}^{\text{OC}}$ is the all-in per-FEU rate
        (base + BAF + THC POL + THC POD).
  \item \textit{Inland:} Cost scales by $n_k$ under the assumption that each FEU moves
        on a separate chassis/truck. This is the standard for FCL drayage.
  \item \textit{Dwell:} Port storage cost (beyond free time) is deferred to P1. In MVP,
        dwell time enters only as a feasibility parameter, not a cost term.
\end{itemize}

%%--------------------------------------------------------------------
\section{Constraints}
%%--------------------------------------------------------------------

\subsection{Flow Conservation}

For each commodity $k \in K$ and each node $n \in N$:

\begin{equation}
  \sum_{\substack{(i,n) \in A^k}} x_{in}^k \;-\; \sum_{\substack{(n,j) \in A^k}} x_{nj}^k \;=\; b_k(n)
  \qquad \forall\, k \in K,\; n \in N
  \label{eq:flow}
\end{equation}

where
\[
  b_k(n) = \begin{cases}
    +1 & \text{if } n = o(k) \quad \text{(commodity departs origin)} \\
    -1 & \text{if } n = d(k) \quad \text{(commodity arrives at destination)} \\
    \phantom{+}0 & \text{otherwise}
  \end{cases}
\]

Flow conservation ensures each commodity is routed along exactly one complete
end-to-end path: one pre-carriage arc, one ocean sailing, the dwell arc at the
selected POD, and one inland arc. With the feasibility pre-filter applied, any
feasible solution automatically satisfies time-window requirements.

\subsection{Ocean Sailing Capacity}

For each ocean sailing arc $(j, p) \in A_{\text{OC}}$:

\begin{equation}
  \sum_{k \in K} n_k\, x_{jp}^k \;\leq\; \text{cap}_{jp}
  \qquad \forall\, (j,p) \in A_{\text{OC}}
  \label{eq:cap}
\end{equation}

The total FEU demand on any single sailing cannot exceed the carrier-allocated slot
capacity. This constraint \emph{couples} routing decisions across commodities: two
commodities from different origins may share the same sailing, and their container
counts jointly consume the finite slot pool. This is the core multi-commodity
coupling constraint.

\subsection{Budget Cap (Optional)}

For each commodity $k \in K$ with finite $B_k$:

\begin{equation}
  \sum_{(i,j) \in A^k_{\text{PC}}} c_{ij}^{\text{PC}}\, x_{ij}^k
  + \sum_{(j,p) \in A^k_{\text{OC}}} n_k\, c_{jp}^{\text{OC}}\, x_{jp}^k
  + \sum_{(p_E,h) \in A^k_{\text{IL}}} n_k\, c_{p_E h}^{\text{IL}}\, x_{p_E h}^k
  \;\leq\; B_k
  \label{eq:budget}
\end{equation}

\subsection{Variable Domains}

\begin{equation}
  x_{ij}^k \in \{0, 1\} \qquad \forall\, k \in K,\; (i,j) \in A^k
  \label{eq:domain}
\end{equation}

%%--------------------------------------------------------------------
\section{Complete Formulation}
%%--------------------------------------------------------------------

\begin{align}
  \min \quad & Z = \sum_{k \in K} \left[
    \sum_{(i,j) \in A^k_{\text{PC}}} c_{ij}^{\text{PC}} x_{ij}^k
    + \sum_{(j,p) \in A^k_{\text{OC}}} n_k c_{jp}^{\text{OC}} x_{jp}^k
    + \sum_{(p_E,h) \in A^k_{\text{IL}}} n_k c_{p_E h}^{\text{IL}} x_{p_E h}^k
  \right] \tag{P} \label{eq:P} \\[4pt]
  \text{s.t.} \quad
  & \sum_{(i,n) \in A^k} x_{in}^k - \sum_{(n,j) \in A^k} x_{nj}^k = b_k(n)
    && \forall\, k \in K,\; n \in N \tag{P.1} \\[2pt]
  & \sum_{k \in K} n_k x_{jp}^k \leq \text{cap}_{jp}
    && \forall\, (j,p) \in A_{\text{OC}} \tag{P.2} \\[2pt]
  & \text{(P.3): budget constraints \eqref{eq:budget} for commodities with finite } B_k \notag \\[2pt]
  & x_{ij}^k \in \{0,1\}
    && \forall\, k \in K,\; (i,j) \in A^k \tag{P.4}
\end{align}

\begin{remark}[Problem structure]
Problem (P) is a Binary Multi-Commodity Network Flow (BMCNF). The LP relaxation
(replacing binary constraints with $x_{ij}^k \in [0,1]$) is typically tight for
network flow problems (total unimodularity in pure flow networks). However, the
capacity constraints \eqref{eq:cap} couple commodities and break total unimodularity,
making integrality of the LP relaxation non-guaranteed in general. In practice, for
small to medium instances, HiGHS will solve the MIP to optimality quickly.
\end{remark}

\begin{remark}[Infeasibility]
A commodity $k$ may be infeasible if no path exists that meets its deadline (empty
$A^k$ after pre-filtering). The solver should detect this during pre-filtering and
report it as an infeasible demand rather than passing an infeasible MILP. In that
case, the agent should surface the infeasibility with a recommended recovery action
(e.g., relax the deadline, consider air mode).
\end{remark}

%%--------------------------------------------------------------------
\section{Transit Time Estimation}
\label{sec:transit}
%%--------------------------------------------------------------------

Arc transit times are used in feasibility pre-filtering and in the instance generator.
They are not optimization variables in MVP; in P1, they become the parameters of
the probabilistic transit time distributions.

\subsection{Trucking (Pre-Carriage and Inland)}

Given origin coordinates $(\phi_i, \lambda_i)$ and destination coordinates $(\phi_j, \lambda_j)$:

\paragraph{Step 1 --- Geodesic distance (Haversine).}
\[
  d_{ij}^{\text{geo}} = 2R \arcsin\!\left(\sqrt{
    \sin^2\!\tfrac{\phi_j - \phi_i}{2}
    + \cos\phi_i \cos\phi_j \sin^2\!\tfrac{\lambda_j - \lambda_i}{2}
  }\right), \qquad R = 6{,}371\,\text{km}
\]

\paragraph{Step 2 --- Road distance correction.}
\[
  \text{dist}_{ij}^{\text{road}} = \rho \cdot d_{ij}^{\text{geo}}
\]
where $\rho$ is a region-specific road factor:
\begin{itemize}[noitemsep]
  \item China pre-carriage (Pearl River Delta, Yangtze Delta): $\rho = 1.25$
  \item US inland (OTR, point-to-point): $\rho = 1.20$
\end{itemize}

\paragraph{Step 3 --- Transit time.}
\[
  \mu_{ij}^{\text{truck}} = \frac{\text{dist}_{ij}^{\text{road}}}{v^{\text{truck}}}
\]
\[
  v^{\text{truck}} = \begin{cases}
    800\,\text{km/day} & \text{US OTR (accounts for HOS 11-hour driving limit)} \\
    600\,\text{km/day} & \text{China pre-carriage (congestion, regulations)}
  \end{cases}
\]

\paragraph{Step 4 --- Standard deviation.}
\[
  \sigma_{ij}^{\text{truck}} = 0.15 \cdot \mu_{ij}^{\text{truck}} \qquad \text{(15\% CV, empirical for road freight)}
\]

\subsection{Ocean Sailing}

\paragraph{Step 1 --- Geodesic distance.} Same Haversine formula.

\paragraph{Step 2 --- Sailing distance.} Ocean routing deviates from great-circle due
to landmasses and shipping lane waypoints.
\[
  \text{dist}_{jp}^{\text{sail}} = \rho^{\text{sail}} \cdot d_{jp}^{\text{geo}},
  \qquad \rho^{\text{sail}} = 1.15 \text{ (Trans-Pacific)}
\]

\paragraph{Step 3 --- Transit time.}
\[
  \mu_{jp}^{\text{OC}} = \frac{\text{dist}_{jp}^{\text{sail}}}{v^{\text{vessel}}},
  \qquad v^{\text{vessel}} = 400\,\text{km/day} \;\approx\; 9\,\text{knots}
\]
(Average container carrier service speed, inclusive of port approach and anchorage wait.)

\paragraph{Step 4 --- Standard deviation.}
\[
  \sigma_{jp}^{\text{OC}} = 0.10 \cdot \mu_{jp}^{\text{OC}} \qquad \text{(10\% CV, empirical for liner services)}
\]

\begin{remark}
These are parametric estimates for instance generation and MVP feasibility checks.
Production models will fit lane-carrier-season-specific distributions from historical
AIS data (NOAA for US waters; commercial AIS for full voyage coverage).
\end{remark}

%%--------------------------------------------------------------------
\section{Instance Generator Specification}
\label{sec:instance}
%%--------------------------------------------------------------------

The instance generator produces synthetic but geographically and commercially
realistic problem instances for solver testing and validation. Building the generator
is a prerequisite to any solver benchmarking.

\paragraph{Note on scope.} The instance generator is to be \emph{reviewed with the developer
before coding begins}. This specification defines the interface and algorithm; the
implementation is a joint session.

\subsection{Public Data Sources}

\begin{center}
\begin{tabular}{lll}
\toprule
Data & Source & Use \\
\midrule
Port lat/lon and UN/LOCODE & UN/LOCODE database (free) & POL and POD node coordinates \\
City lat/lon & GeoNames (free) / OpenStreetMap Nominatim & Origin and destination geocoding \\
Ocean sailing ETDs & Carrier websites (public) & Sailing schedule arc construction \\
Trucking transit time & Haversine + road factor (Section~\ref{sec:transit}) & Pre-carriage and inland arcs \\
Rate parameters & Synthetic distributions (BTS FAF for lane structure) & Ocean rates, inland costs \\
\bottomrule
\end{tabular}
\end{center}

\subsection{Generator Configuration Parameters}

\begin{center}
\begin{tabular}{ll}
\toprule
Parameter & Description \\
\midrule
$|K|$ & Number of commodities to generate \\
$\mathcal{O}$ & Origin cities: Shenzhen, Guangzhou, Suzhou, Shanghai, Ningbo \\
$\mathcal{D}$ & Destination cities: Chicago, Dallas, Los Angeles, New York, Atlanta, Seattle \\
$[v_{\min}, v_{\max}]$ & Volume range, e.g.\ $[10, 300]$ CBM \\
$[w_{\min}, w_{\max}]$ & Weight range, e.g.\ $[500, 20\,000]$ kg \\
Cargo-ready spread & Uniform$[0, 14]$ days from scenario start \\
Deadline spread & Uniform$[30, 60]$ days from cargo-ready date \\
\midrule
\multicolumn{2}{l}{\textit{Synthetic rate parameters}} \\
Pre-carriage rate & \$150--800 (function of road distance) \\
FEU base rate & \$2\,000--6\,000 (Trans-Pacific, lane-specific) \\
BAF & 15--25\% of base rate \\
THC POL & \$150--300/FEU \\
THC POD & \$300--500/FEU \\
Sailing capacity & Uniform$[50, 200]$ FEU per sailing \\
Inland cost (truck) & \$1\,500--6\,000 (function of road distance) \\
\bottomrule
\end{tabular}
\end{center}

\subsection{Generator Algorithm}

\begin{enumerate}
  \item Load POL and POD node coordinates from UN/LOCODE for the ports in scope
        (SHA, NGB, SZX on China side; USLAX, USLGB, USSEA on US side).
  \item For each commodity $k = 1, \ldots, |K|$:
    \begin{enumerate}
      \item Sample origin city from $\mathcal{O}$; geocode to lat/lon.
      \item Sample destination city from $\mathcal{D}$; geocode to lat/lon.
      \item Sample $v_k \sim \text{Uniform}(v_{\min}, v_{\max})$,
            $w_k \sim \text{Uniform}(w_{\min}, w_{\max})$.
      \item Compute $n_k = \ceil{v_k / 67}$.
      \item Sample $t_k^{\text{rdy}}$ (cargo-ready date); set $T_k^{\text{dead}} = t_k^{\text{rdy}} + \text{Uniform}(30, 60)\,\text{d}$.
    \end{enumerate}
  \item Construct $A_{\text{PC}}$: for each (origin city, POL) pair, compute
        $(\text{dist}, \mu, \sigma)$ via Section~\ref{sec:transit}; sample cost
        from rate distribution.
  \item Construct $A_{\text{OC}}$: load or synthesize a weekly sailing schedule for
        each (POL, POD) pair; assign ETD, CYC (ETD $-$ 4 days), carrier, vessel name;
        compute $(\mu, \sigma)$ via Section~\ref{sec:transit}; sample FEU rate and
        capacity.
  \item Construct $A_{\text{DW}}$: one arc per POD with port-specific constants from
        Section 4.4.
  \item Construct $A_{\text{IL}}$: for each (POD exit, destination city) pair, compute
        $(\text{dist}, \mu, \sigma)$ via Section~\ref{sec:transit}; sample cost from
        rate distribution.
  \item Run feasibility pre-filter (Section~\ref{sec:prefilter}); mark infeasible
        commodities for reporting.
  \item Serialize instance to JSON with schema matching the parameter tables above.
\end{enumerate}

%%--------------------------------------------------------------------
\section{Deferred: P1 Extensions}
\label{sec:deferred}
%%--------------------------------------------------------------------

The following items are explicitly out of MVP scope. Each is tracked as a P1 item in
the PRD.

\subsection{Commodity-Specific Customs Inspection Model (P1)}
\label{sec:p1-customs}

In MVP, dwell time is a port-level constant. In P1, it becomes a function of the
commodity and the importer:

\[
  \mu_k^{\text{DW}}(p) = \mu_p^{\text{unload}} + t^{\text{no-exam}} + p_k^{\text{exam}}(p) \cdot t^{\text{exam}}
\]

where the inspection probability is:

\[
  p_k^{\text{exam}}(p) = \operatorname{clip}\!\left(
    \lambda_p^{\text{base}} \cdot \rho_k^{\text{HS}} \cdot \rho_k^{\text{orig}} \cdot
    (1 - \delta_k^{\text{CTPAT}}),\; 0,\; 1
  \right)
\]

\begin{center}
\begin{tabular}{lll}
\toprule
Parameter & Description & Source \\
\midrule
$\lambda_p^{\text{base}}$ & Baseline random examination rate at port $p$ (2--5\%) & CBP statistics (public) \\
$\rho_k^{\text{HS}}$ & HS chapter risk multiplier: 1.0 (low), 2.0 (medium), 4.0 (high) & HS code of commodity \\
$\rho_k^{\text{orig}}$ & Country-of-origin risk multiplier & Country of origin \\
$\delta_k^{\text{CTPAT}}$ & C-TPAT discount: 0 (not certified), 0.50 (certified partner) & Importer attribute \\
\bottomrule
\end{tabular}
\end{center}

\textbf{HS risk tier mapping (examples):}
\begin{itemize}[noitemsep]
  \item Low (1.0): industrial machinery Ch.\ 84, electronics Ch.\ 85 from low-risk origins
  \item Medium (2.0): textiles Ch.\ 50--63, footwear Ch.\ 64
  \item High (4.0): food and agricultural Ch.\ 2--21, chemicals Ch.\ 28--29, certain
        electronics from flagged origins
\end{itemize}

\subsection{Probabilistic Delivery Constraint (P1)}

Replace deterministic feasibility check with a chance constraint:
\[
  \Pr\!\left(\text{delivery time}_k \leq T_k^{\text{dead}}\right) \geq \alpha_k
\]
where $\alpha_k \in \{0.80, 0.90, 0.95\}$ maps to service levels (Economy, Standard, Express).

Since delivery time is a sum of correlated arc-level random variables, this requires
one of: (a) scenario expansion, (b) distributionally robust reformulation, or
(c) Monte Carlo feasibility check on top of the deterministic MILP. Trade-offs to be
evaluated in P1 design review.

\subsection{Multi-Objective Pareto Frontier (P1)}

Extend the objective to return the Pareto frontier of (cost, expected transit time,
on-time probability) using an epsilon-constraint or weighted-sum method.

\subsection{Spot Rate Distributions (P1)}

Ocean arc rates are currently deterministic parameters. In P1, model rates as
lane-specific distributions and minimize expected cost.

\subsection{TEU Option (P1)}

Allow the model to select between FEU and TEU per commodity based on cargo density
(high weight-to-volume ratio may favour a 20$'$ container). Add integer variable
for container type selection.

\subsection{Port Storage Cost (P1)}

Add per-diem storage cost $s_p$ for containers that exceed free-time $f_p$ at the POD.
Requires tracking expected discharge-to-pickup time relative to $f_p$.

%%--------------------------------------------------------------------
\section{Open Items}
%%--------------------------------------------------------------------

\begin{enumerate}
  \item \textbf{Inland cost multi-FEU:} Confirmed --- inland arc cost scales by $n_k$
        (one chassis per FEU). Update \eqref{eq:obj} accordingly (already reflected).

  \item \textbf{Carrier allocation caps:} The PRD mentions carrier allocation caps as a
        constraint (tool \texttt{carrier\_select}). This formulation does not yet model
        allocation caps (monthly contracted min/max volume per carrier per lane). Add
        as a constraint in P1 if needed for the MVP tool set.

  \item \textbf{Infeasible commodity handling:} When no feasible path exists for commodity
        $k$, the solver should return a structured infeasibility report, not just an
        infeasible MILP flag. Define the infeasibility report schema before coding.

  \item \textbf{Multiple sailings per week:} The instance generator should support multiple
        sailings per (POL, POD) pair per week, as is realistic. Confirm arc generation
        logic handles this without duplicating capacity.

  \item \textbf{Transshipment:} Not modeled in MVP. Some Trans-Pacific lanes involve a
        feeder vessel + mainliner (e.g., via Singapore or Kaohsiung). Add transshipment
        arcs in P1.
\end{enumerate}

\end{document}
